\scanpage{57}
\mspage{1}
% archive sheet number: 56
Le dérivateur HOT
\margin{14.11.90\\ (Les premières feuilles, d'il y a deux semaines, ont été depuis, dépassées par \uline{VII})}

\deleted{Il faut définir}
Je définis $\mathrm{Hot}$ à partir de la catégorie de
modèles
\[
M=\mathrm{Cat},\qquad W\subset\mathrm{Fl}(\mathrm{Cat}),
\]
$W$ étant en principe les $\mathrm{Hot}$-équivalences. Mais je
vais dans la mesure du possible axiomatiser plus,
en supposant donné un ensemble $W$ plus
général, appelé ``localisateur fondamental'',
(\unsure{basic localizer|basic localiser}), p.ex. $W=\deleted{\unsure{\circ}}W_{\unsure{\mathbb D|\mathscr D}}$, où $\mathbb D$ est un
dérivateur déjà donné (mais en supposant que
son domaine soit $(\mathrm{Cat})$ tout entier). Nous
allons noter au fur et à mesure de quelles
propriétés de $W$ nous aurons besoin. Aux fins de
référence ultérieure, je vais rappeler ici celles que
j'ai dégagées dans P.S.

\uline{L 1}\quad Saturation de $W$ \deleted{\unsure{$\mathbb D$}}
\[
\left\{\begin{aligned}
&\text{a) les identités dans }W\\
&\text{b) Si deux parmi }gf,f,g\text{ dans }W,\text{ le troisième aussi}
\end{aligned}\right.
\]
\diagram{\begin{tikzcd}[ampersand replacement=\&,column sep=small]
X\arrow[r,"f",dash]\&Y\arrow[r,"g",dash]\&Z
\end{tikzcd}}
c) Si $gf$ et $fg$ sont dans $W$, alors $f$ et $g$ aussi
\diagram{\begin{tikzcd}[ampersand replacement=\&]
X\arrow[r,shift left,"f"]\&Y\arrow[l,shift left,"g"]
\end{tikzcd}}

\diagram{\begin{tikzcd}[ampersand replacement=\&,column sep=small,row sep=small]
\&\text{\uline{L 2}}\arrow[d,Rightarrow]\&
\parbox[t]{.64\linewidth}{$\Delta^1$ est univ. $W$-asph. i.e. $\Delta^1\times X\mathrel{\unsure{-|\to}}X$ est dans $W$ pour tout $X$ dans $\mathrm{Cat}$. Cela implique}\\
\&\text{\uline{L 2 bis}}\&
\parbox[t]{.64\linewidth}{$\Delta^1$ est $W$-\unsure{asph.} (C'était L3$'$ dans P.S.)}\\
\&\text{\uline{L 3}}\&
\parbox[t]{.64\linewidth}{Si $X$ a un objet final, $X$ est $W$-asphérique.}\\
\&\&\parbox[t]{.64\linewidth}{Je note ici variantes renforcées\quad $\big|$ Il y a un L3$'$ dual de L3}\\
\&\text{\uline{L 3} bis}\arrow[uu,Rightarrow]\&
\parbox[t]{.64\linewidth}{Si $X$ est $\mathrm{Hot}$-asphérique, il est $W$-asphérique}\\
\&\text{\uline{L 3} ter}\arrow[u,Rightarrow]\arrow[uuuuu,Rightarrow,bend left=48]\arrow[dl]\&
\parbox[t]{.64\linewidth}{Si $X$ est $\mathrm{Hot}$-asphérique, il est universellement $W$-asphérique, i.e $X\times Y\xrightarrow{\mathrm{pr}_2}Y$ est dans $W$ pour tout $Y$ dans $\mathrm{Cat}$.}\\
\parbox{.13\linewidth}{suffit pour entraîner L2,L3}\&\text{\uline{L 3} quater}\arrow[u,Rightarrow]\&
\parbox[t]{.64\linewidth}{Toute $\mathrm{Hot}$-équivalence est dans $W$}
\end{tikzcd}}
